\documentclass[12pt,a4paper]{article}
\usepackage{fontspec}
\setmainfont{TeX Gyre Pagella}
\usepackage{microtype}
\usepackage[margin=2.5cm]{geometry}
\usepackage{amsmath,amssymb,amsthm}
\usepackage[colorlinks=true,linkcolor=blue!60!black,citecolor=blue!60!black,urlcolor=blue!60!black]{hyperref}
\usepackage{setspace}
\usepackage{titlesec}
\usepackage{enumitem}
\setstretch{1.15}
\titleformat{\section}{\Large\bfseries}{\thesection}{1em}{}
\titleformat{\subsection}{\large\bfseries}{\thesubsection}{1em}{}
\setlength{\parskip}{0.5em}
\setlength{\parindent}{0em}
\newtheorem{theorem}{Theorem}
\newtheorem{prediction}[theorem]{Prediction}

\begin{document}

\begin{center}
{\LARGE\bfseries Wave Function Collapse as Irreversible Entropy Production}\\[18pt]
{\large Paper~II of the Relational Actualism Series}\\[12pt]
{\large Joshua F.\ Sandeman}\\[4pt]
{\small Independent Researcher $\cdot$ Salem, Oregon}\\
{\small\texttt{sansuikyo@gmail.com} $\cdot$ \texttt{relationalactualism.org}}\\[6pt]
{\small April 2026}
\end{center}
\vspace{1em}

\textbf{Wave Function Collapse as Irreversible Entropy Production}

\emph{Paper II of the Relational Actualism Series}

Joshua F. Sandeman

Independent Researcher · Salem, Oregon

sansuikyo@gmail.com · relationalactualism.org

April 2026

\textbf{Abstract}

The quantum measurement problem --- how and when a quantum superposition
becomes a definite classical outcome --- has resisted resolution for
nearly a century. We show that if physical reality is an irreversibly
growing directed acyclic graph governed by the unique self-consistent
growth rule established in Paper I, the measurement problem does not
arise. An interaction becomes a permanent event in the causal graph when
it produces an irreversible increase in quantum relative entropy ΔS(ρ ‖
σ\textsubscript{0}) above a specific threshold ΔS* = 0.601 nats,
computed from the BDG action at the Planck density with no free
parameters. Below this threshold, interactions are virtual (quantum,
reversible, off-shell). Above it, they are actual (classical,
irreversible, on-shell). There is no Heisenberg cut, no
observer-dependent collapse, and no modification of quantum mechanics
--- only an objective, observer-independent criterion for when events
become real. The criterion is Lorentz-covariant (proved), causally
invariant (Lean 4 verified), and consistent with the Born rule. We
derive three specific, falsifiable predictions: (i) a strict null result
for gravity-mediated entanglement experiments (the BMV protocol); (ii) a
parameter-free spin-bath collapse timescale t* = 0.274/g; and (iii) a
hard ceiling on fault-tolerant quantum array size N\textsubscript{max} =
η × p\textsubscript{th} (the Kinematic Coherence Bound). Each prediction
distinguishes this framework from all existing interpretations of
quantum mechanics.

\textbf{1. Introduction}

The measurement problem is the central conceptual puzzle of quantum
mechanics. The Schrödinger equation describes a world of superpositions:
a particle can be in two places at once, a cat can be simultaneously
alive and dead, the moon's position is indefinite until observed. Yet
every observation, without exception, finds a definite, particular
outcome. How does the quantum realm of possibilities become the
classical realm of facts?

A century of theoretical effort has produced no consensus. The
Copenhagen interpretation draws a line --- the Heisenberg cut ---
between the quantum system and the measuring apparatus, and declines to
say what happens at the line. Many-worlds posits that every outcome
occurs in a branching multiverse, eliminating collapse at the cost of an
enormous ontological commitment. Spontaneous collapse models (GRW, CSL)
modify the Schrödinger equation with stochastic terms, introducing new
parameters. Pilot-wave theories (de Broglie-Bohm) add hidden variables.
Each approach either adds machinery, changes the mathematics, or denies
the problem exists.

This paper takes a different approach, following from the hypothesis
established in Paper I {[}1{]}: physical reality is an irreversibly
growing directed acyclic graph (DAG) whose growth rule is uniquely
determined by self-consistency. If this hypothesis is correct, the
measurement problem \emph{does not arise}, because the framework
contains an objective, observer-independent criterion for when quantum
possibilities become classical facts. The criterion is not imposed ---
it is derived from the unique growth rule. And it makes specific,
falsifiable predictions that distinguish it from every existing
interpretation.

The paper is organised as follows. Section 2 derives the actualization
criterion from the BDG action. Section 3 shows how this dissolves the
measurement problem. Section 4 proves Lorentz covariance. Section 5
resolves the Unruh paradox. Section 6 presents the causal invariance
theorem. Section 7 discusses Born rule consistency. Section 8 derives
the BMV null prediction. Section 9 derives the spin-bath collapse
timescale. Section 10 presents the Kinematic Coherence Bound for quantum
computing. Section 11 compares with other approaches. Section 12
concludes.

\textbf{2. The Actualization Criterion}

\textbf{2.1 On-shell and off-shell in the causal graph}

In the growing DAG, each vertex represents an actualization event --- an
irreversible physical interaction that writes a permanent edge into the
causal graph. Not every proposed interaction becomes a vertex. The BDG
growth rule (Paper I, Section 2) assigns an action S to each proposed
vertex based on its causal neighbourhood. Vertices with S \textgreater{}
0 are accepted; those with S ≤ 0 are rejected.

The physical distinction between accepted and rejected vertices
corresponds to the quantum field theory distinction between on-shell and
off-shell processes. An on-shell process --- one satisfying the
mass-shell condition p\textsuperscript{μ}p\textsubscript{μ} = m²c² ---
produces a real particle that propagates to a detector. An off-shell
process --- a virtual particle exchange --- contributes to scattering
amplitudes but does not produce a detectable particle. In RA, this
distinction is not a matter of mathematical convenience but of
ontological status: on-shell processes are vertices in the graph;
off-shell processes are not.

\textbf{2.2 The entropic formulation}

The actualization criterion can be stated in terms of quantum relative
entropy. The quantum relative entropy S(ρ ‖ σ\textsubscript{0}) measures
how distinguishable the state ρ (produced by the interaction) is from
the vacuum state σ\textsubscript{0}. An interaction actualises when:

\begin{quote}
\emph{ΔS(ρ ‖ σ₀) \textgreater{} ΔS*}
\end{quote}

where ΔS* = 0.601 nats is the actualization threshold. This threshold is
computed from the BDG action at the Planck density μ = 1 (Paper I,
Section 2): ΔS* = −log P\textsubscript{acc}(1) = −log(0.548) = 0.601
nats {[}CV{]}. It is the unique self-consistent growth rule's
selectivity at its natural operating point.

The entropic formulation is primary; the on-shell condition
p\textsuperscript{μ}p\textsubscript{μ} = m²c² is its perturbative-regime
implementation. In the perturbative regime, the two criteria coincide:
an on-shell interaction produces a state sufficiently distinguishable
from the vacuum to exceed ΔS*. Off-shell interactions have ΔS = 0 (they
do not change the distinguishability from the vacuum) and therefore do
not actualise.

\textbf{2.3 The threshold is not a free parameter}

The uniqueness theorem of Paper I establishes that the BDG growth rule
is the only self-consistent option. Therefore ΔS* = 0.601 nats is not
one threshold among many --- it is the only threshold compatible with
the self-consistency of the causal graph. No alternative actualization
criterion exists within the framework.

\textbf{3. The Measurement Problem Dissolved}

\textbf{3.1 No Heisenberg cut}

In RA, a quantum state is any state with ΔS(ρ ‖ σ\textsubscript{0})
below the actualization threshold ΔS*. It has not yet been written into
the causal graph as a permanent event. An actualization event occurs
when, through interaction with an environment, the relative entropy
crosses the threshold. A new permanent vertex is written. Nothing
``collapses'' --- the wavefunction does not undergo a sudden,
discontinuous change. What happens is that the causal graph gains a new
vertex, irreversibly.

There is no Heisenberg cut because there is no fundamental distinction
between quantum and classical. There is only a gradient of actualization
density. The Schrödinger equation is the Level 3 (macroscopic)
approximation that applies when the actualization density is low: many
proposed events, few actualising. Classical mechanics is the Level 3
approximation that applies when the density is high: most proposed
events actualise, behaviour is deterministic. Both are approximations to
the same underlying discrete process.

\textbf{3.2 Schrödinger's cat}

Schrödinger's cat is not in a superposition of alive and dead states. A
real cat is a densely self-coupled actualization network --- trillions
of atoms exchanging on-shell bosons millions of times per second. The
actualization density is astronomical. The cat is classical by
constitution, not by observation. The radioactive isotope alone is the
genuine locus of quantum indeterminacy: a single system waiting for its
first actualization event. The paradox dissolves because the thought
experiment misidentified the quantum system. The cat was never quantum;
only the isotope was.

\textbf{3.3 The vacuum does not gravitate}

A direct consequence of the actualization criterion: off-shell processes
have ΔS = 0 and do not write vertices into the graph. Since the
gravitational field is sourced only by actualized vertices (via the
actualization projector P\textsubscript{act} applied to the
stress-energy tensor), virtual quantum fluctuations do not gravitate.
The cosmological constant is not small --- it is identically zero,
because the vacuum (the state of all-unactualized-candidates)
contributes nothing to the metric source. This dissolves the
cosmological constant problem structurally, as a consequence of the same
criterion that dissolves the measurement problem.

\textbf{4. Lorentz Covariance}

The actualization criterion must be frame-independent: two observers in
relative motion must agree on whether an event has actualised. This is
guaranteed by the frame-independence of quantum relative entropy.

\textbf{Theorem (Frame Independence) {[}LV{]}.} \emph{For any unitary U
satisfying Uσ}\textsubscript{0}\emph{U† = σ}\textsubscript{0}\emph{:
S(UρU† ‖ σ}\textsubscript{0}\emph{) = S(ρ ‖ σ}\textsubscript{0}\emph{).}

This is proved in Lean 4 (frame\_independence in
RA\_AQFT\_Proofs\_v10.lean). For Poincaré transformations U, the
Minkowski vacuum σ\textsubscript{0} is invariant (Uσ\textsubscript{0}U†
= σ\textsubscript{0}), so the theorem applies. The actualization
criterion is therefore Lorentz-covariant: if ΔS \textgreater{} ΔS* in
one frame, it exceeds ΔS* in all frames.

This resolves a conceptual difficulty that afflicts many collapse
models: the question of which frame defines the collapse surface. In RA,
no such question arises. The criterion is frame-independent by the
unitary invariance of relative entropy.

\textbf{5. The Unruh Resolution}

The Unruh effect {[}2{]} poses a challenge to any observer-independent
collapse criterion: a uniformly accelerated (Rindler) observer perceives
the Minkowski vacuum as a thermal bath at the Unruh temperature
T\textsubscript{U} = ħa/(2πck\textsubscript{B}). If thermal excitations
trigger actualization, the Rindler observer would see events that the
inertial observer does not, violating observer-independence.

\textbf{Theorem (Rindler Stationarity) {[}LV{]}.} \emph{d/dt
S(α}\textsubscript{t}\emph{(ρ}\textsubscript{β}\emph{) ‖
σ}\textsubscript{0}\emph{) = 0 under the modular flow
α}\textsubscript{t} \emph{generated by the Rindler Hamiltonian.}

This is proved in Lean 4 (rindler\_stationarity in
RA\_AQFT\_Proofs\_v10.lean). The KMS thermal state ρ\textsubscript{β}
associated with the Rindler wedge has identically zero time-derivative
of relative entropy under the modular flow. The threshold ΔS* is never
crossed. No actualization occurs.

The physical interpretation: the Rindler thermal bath is a
\emph{relabelling} of the Minkowski vacuum, not a new physical state.
The relative entropy between the relabelled state and the vacuum is
constant --- the state is not becoming more distinguishable from the
vacuum over time. Both the inertial and accelerated observers agree: no
physical event has occurred. The Unruh ``particles'' are a coordinate
artefact, not actualization events.

\textbf{6. Causal Invariance}

The quantum measure for a set of spacelike-separated actualization
events must be independent of the order in which they are processed.
This is the causal invariance requirement, essential for consistency of
the growth dynamics.

\textbf{Theorem (Amplitude Locality) {[}LV{]}.} \emph{For the BDG
amplitude function, a(v \textbar{} C) = a(v \textbar{} C') whenever C ∩
past(v) = C' ∩ past(v).}

This is proved in Lean 4 (bdg\_amplitude\_locality in
RA\_AmpLocality.lean) as a theorem of the discrete DAG dynamics, without
any appeal to QFT microcausality. The proof exploits the transitivity of
the causal order: causal intervals {[}u, v{]} lie entirely within
past(v), so the BDG action increment at v depends only on the portion of
the causal set within v's past.

\textbf{Theorem (Causal Invariance) {[}LV, conditional{]}.} \emph{Given
amplitude locality, the quantum measure μ}\textsubscript{σ}\emph{(S) =
μ}\textsubscript{σ'}\emph{(S) for any two linear extensions σ, σ' of the
causal order.}

With amplitude locality now proved as a theorem (rather than an axiom),
causal invariance is established unconditionally for the BDG dynamics.
The one remaining sorry in the Lean formalisation is a List.Perm
induction step --- a mechanical combinatorial fact about commuting
complex multiplication, not a conceptual gap.

\textbf{7. Born Rule Consistency}

The Born rule --- that the probability of a measurement outcome is the
squared modulus of the wavefunction --- is consistent with RA's quantum
measure but is not \emph{derived} from it. In RA, the Born rule is a
\emph{constraint} on the amplitude function: a(x) ∝ ψ(x). The quantum
measure, which generalises Sorkin's classical quantum measure framework
{[}3{]}, assigns probabilities to sets of actualization events. These
probabilities are consistent with the Born rule when the amplitude
function is identified with the wavefunction.

This is a sharper epistemic claim than ``many-worlds derives the Born
rule'': RA's quantum measure is a specific mathematical structure (the
squared norm of the product of local amplitudes over a causal set), and
the Born rule falls out as the single-event marginal of this structure.
It is not independently postulated; it is a consequence of the
multiplicative amplitude framework. But it is a consequence of the
\emph{structure} of the quantum measure, not of the dynamics --- which
is why we call it ``consistency'' rather than ``derivation.''

\textbf{8. The BMV Null Prediction}

\textbf{8.1 The experimental protocol}

The Bose-Marletto-Vedral (BMV) protocol {[}4,5{]} places two masses in
quantum spatial superposition and asks whether they can become
gravitationally entangled. If gravity is fundamentally quantum, the two
masses should develop entanglement mediated by a graviton or
gravitational field in superposition. If gravity is fundamentally
classical, no entanglement should develop.

\textbf{8.2 RA's prediction}

In RA, the spacetime metric is sourced by the actualization projector
P\textsubscript{act}{[}T\textsubscript{μν}{]} --- only actualized events
contribute to the gravitational field. A mass in quantum spatial
superposition has not actualised its position (its centre-of-mass state
has ΔS \textless{} ΔS*). Therefore its position is not a vertex in the
causal graph, and it sources zero gravitational field.

The prediction is strict: \textbf{no gravity-mediated entanglement} at
any mass scale where the centre-of-mass can be maintained in
superposition. This is a null result, independent of the mass, the
superposition width, or the duration of the experiment. It holds for the
same reason the vacuum does not gravitate: unactualized degrees of
freedom do not source the metric.

A positive result in any BMV-type experiment would directly falsify RA's
treatment of the gravity-measurement interface. Multiple experimental
groups are currently pursuing this test.

\textbf{9. The Spin-Bath Collapse Timescale}

RA makes a parameter-free prediction for the timescale of wavefunction
actualisation in spin-bath experiments:

\begin{quote}
\emph{t* = (1/g)√(ΔS*/2) = 0.274/g}
\end{quote}

where g is the spin-bath coupling constant and ΔS* = 0.601 nats is the
actualization threshold {[}CV{]}. The derivation: a spin-1/2 system
coupled to a bath of N spins with coupling g accumulates relative
entropy at a rate proportional to g². The threshold ΔS* is crossed at
time t* = (1/g)√(ΔS*/2).

This prediction is specific to RA and distinguishes it from GRW and CSL
collapse models, which predict different functional forms for the
collapse timescale (GRW: exponential in particle number; CSL: dependent
on a free parameter λ). The RA prediction has \emph{no free parameters}:
both g and ΔS* are determined by the physics of the specific
experimental system and the unique BDG growth rule, respectively.

Superconducting circuit experiments operating at millikelvin
temperatures can probe this regime. The prediction is that decoherence
rates will show a threshold behaviour at ΔS = ΔS*, absent in standard
Lindblad master equation predictions.

\textbf{10. The Kinematic Coherence Bound}

Because each qubit is an additional kinematic site at which
actualization events can occur, RA predicts a structural ceiling on the
size of a fault-tolerant quantum array.

\textbf{10.1 The bound}

\textbf{Prediction (KCB).} \emph{The maximum number of logical qubits in
a fault-tolerant quantum array is N}\textsubscript{max} \emph{= η ×
p}\textsubscript{th}\emph{, where η is the single-qubit quality factor
(hardware-dependent) and p}\textsubscript{th} \emph{is the
fault-tolerance threshold (typically 10⁻² to 10⁻³).}

The derivation: each qubit substrate (electron, photon) has a minimum
BDG stability score of 1 {[}LV{]} (structural\_fragility in
RA\_D1\_Proofs.lean). This means each qubit has a nonzero probability
per unit time of undergoing a spurious actualization event --- an
irreversible quantum-to-classical transition that constitutes an error.
For N qubits, the cumulative probability of at least one spurious
actualization scales linearly with N. When this cumulative probability
exceeds the fault-tolerance budget p\textsubscript{th}, the array cannot
maintain coherence regardless of single-qubit fidelity.

\textbf{10.2 Distinguishing prediction}

Standard decoherence theory predicts continuous (though exponentially
difficult) scaling: in principle, with sufficient error correction
overhead, arbitrarily large quantum computers can be built. RA predicts
a \emph{hard wall} at N\textsubscript{max}. The distinction is
qualitative: RA says the scaling terminates; standard theory says it
merely becomes expensive. This is testable at the scale of hundreds to
thousands of logical qubits with surface code experiments.

\textbf{11. Comparison with Other Approaches}

\textbf{11.1 Copenhagen}

Copenhagen draws the Heisenberg cut and declines to specify the collapse
mechanism. RA replaces the cut with an objective criterion (ΔS
\textgreater{} ΔS*) that requires no observer, no measuring apparatus,
and no macroscopic/microscopic distinction.

\textbf{11.2 Many-worlds}

Many-worlds eliminates collapse by positing that every branch of the
wavefunction is equally real. RA eliminates collapse differently: the
wavefunction never collapses because there is no wavefunction in the
fundamental ontology --- there are only actualization events in the
causal graph. The ``wavefunction'' is a Level 3 (macroscopic)
description that applies when the actualization density is low. At the
fundamental level, there is only the growing graph with its threshold.

\textbf{11.3 GRW / CSL}

Spontaneous collapse models modify the Schrödinger equation with
stochastic terms controlled by free parameters (λ, r\textsubscript{c} in
GRW; λ in CSL). RA introduces no free parameters: ΔS* is computed from
the unique BDG action. The collapse timescale in RA (t* = 0.274/g) has a
different functional form from GRW/CSL predictions, providing a clean
experimental discriminant.

\textbf{11.4 de Broglie-Bohm}

Pilot-wave theory adds hidden variables (particle positions) guided by
the wavefunction. RA adds nothing to the ontology --- the causal graph
is the complete description. There are no hidden variables because there
is nothing hidden: the graph is the reality, and its vertices are the
events.

\textbf{11.5 Summary of distinguishing predictions}

Three predictions distinguish RA from all existing interpretations: (i)
the BMV null result (not predicted by many-worlds, pilot-wave, or
standard Copenhagen); (ii) the parameter-free collapse timescale t* =
0.274/g (not predicted by any interpretation without free parameters);
(iii) the hard ceiling N\textsubscript{max} on quantum array size (not
predicted by any interpretation that treats decoherence as continuous
scaling).

\textbf{12. Conclusions}

The measurement problem dissolves within Relational Actualism because
the framework contains an objective, observer-independent criterion for
when quantum possibilities become classical facts: the actualization
threshold ΔS* = 0.601 nats, derived from the unique self-consistent
growth rule of the causal graph (Paper I). The criterion is
Lorentz-covariant (proved), causally invariant (Lean 4 verified), and
consistent with the Born rule. It resolves the Unruh paradox, explains
why the vacuum does not gravitate, and makes three specific, falsifiable
predictions that distinguish it from every existing interpretation of
quantum mechanics.

The DAG hypothesis is posited. The actualization criterion is derived.
The predictions are tests. If the BMV experiment detects
gravity-mediated entanglement, or if quantum computers scale past the
Kinematic Coherence Bound without encountering a hard wall, the
framework is falsified. If they do not, each non-observation adds
evidence to the hypothesis that physical reality is, at its most
fundamental level, an irreversibly growing causal graph.

Companion papers derive further consequences: coupling constants and
particle spectrum (Paper I {[}1{]}), cosmological parameters and
nucleation cosmology (Paper III {[}6{]}), substrate-independent
conditions for biological complexity (Paper IV {[}7{]}), and the
complete framework synthesis (Paper V {[}8{]}). An accessible narrative
overview is given in Paper 0 {[}9{]}.

\textbf{References}

{[}1{]} J. F. Sandeman, Paper I of this series: The Physics of a
Self-Consistent Causal Graph (2026).

{[}2{]} W. G. Unruh, Phys. Rev. D 14, 870 (1976).

{[}3{]} R. D. Sorkin, Mod. Phys. Lett. A 9, 3119--3128 (1994).

{[}4{]} S. Bose et al., Phys. Rev. Lett. 119, 240401 (2017).

{[}5{]} C. Marletto and V. Vedral, Phys. Rev. Lett. 119, 240402 (2017).

{[}6{]} J. F. Sandeman, Paper III of this series (2026). Zenodo:
doi:10.5281/zenodo.19363017.

{[}7{]} J. F. Sandeman, Paper IV of this series (2026).

{[}8{]} J. F. Sandeman, Paper V of this series (2026). Zenodo:
doi:10.5281/zenodo.19363077.

{[}9{]} J. F. Sandeman, Paper 0: A Biography of the Universe (2026).

{[}10{]} D. M. T. Benincasa and F. Dowker, Phys. Rev. Lett. 104, 181301
(2010).

{[}11{]} K. Yeats, Class. Quantum Grav. 42, 145003 (2025).
arXiv:2412.14036.

{[}12{]} G.-C. Rota, Z. Wahrscheinlichkeitstheorie 2, 340--368 (1964).

\end{document}
